\section{Objet et portée de l'annexe}
\label{annexe-J-objet}

Cette annexe formalise la distinction pli/couture introduite dans les analyses préparatoires entre deux
mécanismes de transition observée : la recomposition de l'exécution
hybride et la déformation continue de l'état parcouru à exécution
combinatoirement inchangée. Elle donne un sens précis à la formule :
\begin{quote}
\emph{Ce qui recompose l'exécution recoud la frontière ; ce qui la
déforme la plie.}
\end{quote}

L'énoncé ne repose ni sur l'opposition protocole/état, ni sur une
hypothèse de lissité globale de la frontière. Il s'appuie sur trois
domaines explicitement distingués : le domaine causal des paramètres
admissibles, le domaine markovien sur lequel l'état résumé suffit, et le
domaine réalisable effectivement atteint par les histoires et protocoles
considérés.

\begin{encadre}[title={Résultat directeur}]
Sur chaque cellule où la signature événementielle est constante, la
frontière d'appropriation est la frontière d'une carte de réponse définie
sur une exécution de type combinatoire fixe. Une transition qui reste
dans une telle cellule est une transition de type C : elle plie la
frontière sans recoudre le calendrier. Une transition qui rencontre le
lieu où la signature événementielle change est de type K : elle peut
recoudre des morceaux de frontière issus de cellules différentes. Les
cas d'intersection sont mixtes et ne doivent pas être forcés dans une
dichotomie exclusive.
\end{encadre}

\section{Domaine de travail}
\label{annexe-J-domaines}

\subsection{Domaine causal}

Soit $Q$ un espace de paramètres regroupant, selon la famille étudiée,
un état lent initial, des paramètres de protocole, des durées d'épisode
et des amplitudes d'entrée. Le \emph{domaine causal} $D_{\mathrm c}\subseteq
Q$ contient les paramètres pour lesquels le protocole est défini sans
recourir à une information future. Une coordonnée de $Q$ peut donc être
un paramètre de protocole sans cesser d'être causale, à condition qu'elle
soit fixée avant l'exécution correspondante.

\subsection{Domaine markovien}

Le \emph{domaine markovien} $D_{\mathrm m}\subseteq D_{\mathrm c}$ est le
sous-ensemble sur lequel le résumé déclaré détermine la réponse future à
la sonde. Dans le cadre de l'Appendice~I, il s'agit du domaine où
$z=(p,\chi,\theta,s)$ est comportementalement suffisant. Si une variable
cachée, par exemple $ssc$, devient active sous la loi considérée, le
domaine markovien associé à ce résumé se réduit ou le résumé doit être
agrandi.

\subsection{Domaine réalisable}

Le \emph{domaine réalisable} $D_{\mathrm r}\subseteq D_{\mathrm m}$ est
l'image des histoires et protocoles admissibles effectivement considérés.
Il peut être de dimension plus faible que l'espace formel des états. Les
frontières pertinentes sont toujours prises relativement à
$D_{\mathrm r}$, sauf mention contraire.

On travaillera désormais sur
\[
D:=D_{\mathrm r},
\]
muni de sa topologie relative. Cette convention évite d'attribuer un
statut physique à des points formellement écrits mais inaccessibles.

\section{Signature d'exécution et cellules hybrides}
\label{annexe-J-signature}

Pour $q\in D$, notons
\[
\operatorname{Exec}(q)
\]
l'exécution déterministe produite par le protocole déclaré. On associe à
cette exécution une \emph{signature combinatoire}
\[
\kappa(q)\in\mathcal K,
\]
qui enregistre uniquement les données discrètes choisies : nombre et
type d'événements, ordre des transitions, identité des règles déclenchées
et, si nécessaire, segmentation des épisodes. Les dates exactes et les
états continus aux événements ne font pas partie de $\kappa$.

Le choix de $\kappa$ est relatif à la question étudiée. Une signature
trop grossière peut masquer une recomposition pertinente ; une signature
trop fine peut transformer une simple dérive temporelle en changement
artificiel de calendrier. La déclaration doit donc précéder la
classification.

Pour $k\in\mathcal K$, posons
\[
C_k=\operatorname{int}_D\{q\in D\mid \kappa(q)=k\}.
\]
Les composantes connexes non vides des $C_k$ sont appelées
\emph{cellules hybrides}. Le lieu de recomposition est
\begin{equation}
\label{eq:annexe-J-K}
K_\kappa
=
D\setminus\bigcup_k C_k.
\end{equation}
Il contient les points où la signature n'est pas localement constante.
Il peut comprendre des surfaces de changement d'événement, des coins,
des intersections multiples ou des ensembles plus singuliers.

\begin{encadre}[title={Précaution de portée}]
$K_\kappa$ n'est pas en lui-même une frontière d'appropriation. Il marque
un changement de type d'exécution. Une recomposition du calendrier peut
être observationnellement neutre à la résolution $\rho$, et une bascule
d'issue peut au contraire se produire loin de $K_\kappa$.
\end{encadre}

\section{Carte de réponse par cellule}
\label{annexe-J-reponse-cellule}

Soit
\[
\widehat\Omega_{\rho}:D\longrightarrow\mathcal Y_\rho
\]
la carte de réponse obtenue après fixation de la sonde et du protocole
extérieur. Pour une lecture binaire, on écrit
\[
A=\{q\in D\mid \widehat\Omega_\rho(q)=1\},
\qquad
\Sigma=\partial_D A.
\]

Dans une cellule hybride $C$, la succession des événements est fixe.
Les dates et les états aux coupures peuvent néanmoins varier avec $q$.
Lorsque l'exécution dépend continûment des paramètres à l'intérieur de
$C$, il existe une carte d'état de coupure ou d'observation
\[
E_C:C\longrightarrow \mathcal X_C
\]
et une lecture cellulaire
\[
L_C:\mathcal X_C\longrightarrow\mathcal Y_\rho
\]
telles que
\begin{equation}
\label{eq:annexe-J-factorisation-cellule}
\widehat\Omega_\rho|_C=L_C\circ E_C.
\end{equation}
Cette factorisation ne prétend pas que $E_C$ soit injective ni que
$L_C$ soit continue pour une sortie binaire.

Pour étudier une frontière binaire, on introduit si possible une
observable pré-binaire continue
\[
g_C:C\longrightarrow\mathbb R,
\]
telle que l'issue soit donnée par le signe de $g_C$, hors éventuel
plateau déclaré :
\[
\widehat\Omega_\rho(q)=1
\quad\Longleftrightarrow\quad
g_C(q)>0.
\]

\section{Décomposition de la frontière}
\label{annexe-J-decomposition}

Définissons la partie cellulaire de la frontière par
\[
\Sigma_C=\Sigma\cap C
\]
pour chaque cellule hybride $C$, et la partie de couture par
\[
\Sigma_K=\Sigma\cap K_\kappa.
\]
Comme
\[
D=K_\kappa\sqcup\bigsqcup_C C,
\]
on obtient exactement la décomposition disjointe
\begin{equation}
\label{eq:annexe-J-decomposition}
\boxed{
\Sigma
=
\Sigma_K\sqcup\bigsqcup_C\Sigma_C.
}
\end{equation}
Cette identité est ensembliste et topologique ; elle n'affirme aucune
régularité différentielle. Les adhérences n'interviennent que pour
décrire le raccord éventuel entre des morceaux cellulaires distincts.

\resultstep{prop:J-pli-couture}
\begin{encadre}[title={Proposition \theproposition{} --- Pli et couture de la frontière}]
Supposons :
\begin{enumerate}
\item l'exécution déterministe sur $D$ ;
\item une signature $\kappa$ localement constante hors de $K_\kappa$ ;
\item sur chaque cellule hybride $C$, une factorisation
$\widehat\Omega_\rho|_C=L_C\circ E_C$ ;
\item pour la description locale par niveau, une observable pré-binaire
continue $g_C$ sans plateau de signe au voisinage considéré.
\end{enumerate}
Alors :
\begin{enumerate}
\item dans l'intérieur d'une cellule, la frontière est entièrement
déterminée par la déformation continue de l'exécution de type fixe ;
localement,
\[
\Sigma\cap C=\partial_C\{g_C>0\},
\]
et, si $g_C$ est de classe $C^1$ et zéro une valeur régulière,
\[
\Sigma\cap C=\{g_C=0\};
\]
\item toute transition d'issue le long d'un chemin continu qui reste dans
une même cellule est de type C ;
\item si $\kappa\circ\gamma$ n'est pas localement constante au point de
bascule $t_*$, la transition est de type K ; l'appartenance
$\gamma(t_*)\in K_\kappa$ est nécessaire, mais non suffisante. Si la
transition est simultanément portée par l'adhérence d'un pli cellulaire,
elle est mixte ;
\item les morceaux de frontière provenant de cellules de signatures
différentes ne peuvent être raccordés qu'à travers leur adhérence sur
$K_\kappa$.
\end{enumerate}
\end{encadre}

\paragraph{Démonstration.}
Dans une cellule $C$, $\kappa$ est constante. Toute variation de la
réponse y provient donc de la variation de $E_C(q)$ et de sa lecture par
$L_C$, non d'un changement d'ordre ou de type des événements. Pour une
observable pré-binaire continue, la frontière relative de l'ensemble
positif est contenue dans le niveau zéro ; en l'absence de plateau de
signe, elle lui est égale localement. Si zéro est une valeur régulière,
le théorème des fonctions implicites donne une hypersurface locale.

Considérons ensuite un chemin continu $\gamma:[a,b]\to D$. Si un point
de bascule $t_*$ possède un voisinage envoyé dans une seule cellule, la
signature reste constante autour de $t_*$ : la transition est de type C.
Si $\kappa\circ\gamma$ n'est pas localement constante en $t_*$,
alors nécessairement $\gamma(t_*)\in K_\kappa$ et la transition est de
type K. La réciproque est fausse : un chemin peut toucher ou longer
$K_\kappa$ tout en conservant une signature constante. Si le point est
aussi porté par l'adhérence d'un niveau cellulaire de bascule, les deux
mécanismes coexistent et la transition est mixte.
Enfin, deux cellules distinctes ne se rencontrent pas dans leur intérieur ;
tout raccord entre leurs frontières relatives passe donc par leurs
adhérences, incluses dans le complément des intérieurs, donc sur
$K_\kappa$. \hfill$\square$

\section{Transitions le long d'une famille unidimensionnelle}
\label{annexe-J-chemins}

Soit $\gamma:I\to D$ une famille continue. Un point $t_*$ est une
bascule effective si la lecture binaire prend des valeurs différentes
de part et d'autre de $t_*$. On distingue :

\begin{description}
\item[Type C.] Il existe un voisinage de $t_*$ sur lequel
$\kappa\circ\gamma$ est constante.
\item[Type K.] La signature $\kappa\circ\gamma$ n'est pas localement
constante en $t_*$.
\item[Type mixte.] $\gamma(t_*)$ appartient à la fois au lieu de
recomposition et à l'adhérence d'un pli cellulaire qui change déjà le
signe de l'observable.
\end{description}

Cette classification est locale au chemin et à la signature déclarée.
Elle n'est pas une propriété absolue d'un paramètre. Une même coordonnée
peut produire des transitions C sur une plage et K sur une autre.

\resultstep{cor:J-axe-type}
\begin{encadre}[title={Corollaire \theproposition{} --- La nature de l'axe ne décide pas le type}]
Aucune implication générale ne relie « paramètre de protocole » à type K,
ni « variable d'état » à type C. Une durée de protocole peut agir comme
un paramètre continu de l'état de coupure tant que la signature reste
fixe ; une variable d'état peut au contraire déclencher ou supprimer un
événement et produire une transition K.
\end{encadre}

\section{Durée dans une marche et état de coupure}
\label{annexe-J-duree}

Considérons une durée $T$ variant dans un intervalle où la signature
reste constante. L'exécution avant la coupure est de même type
combinatoire, mais l'état atteint au temps $T$ varie :
\[
X_{\mathrm{cut}}(T)=\Phi_T(X_0).
\]
La réponse postérieure se factorise alors sous la forme
\[
\widehat\Omega_\rho(T)
=
L\bigl(\Psi(X_{\mathrm{cut}}(T))\bigr),
\]
où $\Psi$ désigne l'évolution post-coupure. Tant qu'aucun événement
n'est créé, supprimé ou permuté, $T$ est donc un paramètre d'état de
coupure déguisé. Les transitions observées sont de type C.

À l'extrémité d'une marche, un événement peut apparaître ou disparaître.
Le même axe devient alors capable de produire une transition K. La
classification dépend ainsi du point parcouru, non du nom de la
coordonnée.

\section{Recomposition sans bascule et bascule sans recomposition}
\label{annexe-J-independance}

Les ensembles $K_\kappa$ et $\Sigma$ sont logiquement indépendants.
Quatre situations sont possibles :
\begin{enumerate}
\item $q\notin K_\kappa$ et $q\notin\Sigma$ : exécution et issue sont
localement stables ;
\item $q\notin K_\kappa$ et $q\in\Sigma$ : pli C ;
\item $q\in K_\kappa$ et $q\notin\Sigma$ : recomposition
observationnellement neutre ;
\item $q\in K_\kappa$ et $q\in\Sigma$ : couture K ou transition mixte.
\end{enumerate}

Cette table interdit deux confusions fréquentes : compter les événements
ne suffit pas à prédire l'issue, et détecter une bascule d'issue ne suffit
pas à conclure à une recomposition de l'exécution.

\section{Causalité, suffisance et réalisabilité}
\label{annexe-J-coherence-domaines}

La Proposition~\ref{prop:J-pli-couture} doit être appliquée sur l'intersection des trois
contraintes déclarées.

\paragraph{Causalité.}
Une famille qui ajuste un paramètre en fonction d'un événement futur ne
définit pas un axe causal de protocole. Elle peut servir à une analyse
a posteriori, mais non à une intervention admissible sans reformulation.

\paragraph{Suffisance markovienne.}
La cellule hybride et la carte de coupure doivent être définies dans un
état qui détermine le futur. Si deux histoires de même résumé peuvent
produire des suffixes différents, la frontière projetée peut se
superposer artificiellement. Il faut alors agrandir l'état, par exemple
en ajoutant une variable lente cachée.

\paragraph{Réalisabilité.}
Une couture dans l'espace ambiant n'a pas de conséquence expérimentale
si aucune histoire admissible ne l'atteint. Inversement, une section
réalisable tangente à une frontière peut produire une apparence de
terminaison ou de duplication qui n'existe pas dans l'espace plus grand.

\resultstep{cor:J-domaine-realisable}
\begin{encadre}[title={Corollaire \theproposition{} --- Restriction au domaine réalisable}]
La frontière observable d'une famille est la frontière relative de la
restriction de la carte de réponse au domaine réalisable. Elle n'est
égale à l'intersection avec la frontière ambiante que sous les conditions
de transversalité et de non-coïncidence déjà précisées dans
l'Appendice~I.
\end{encadre}

\section{Structure géométrique et activité décisionnelle}
\label{annexe-J-structure-activite}

Un axe peut déformer fortement l'état de coupure sans changer l'issue à
la résolution considérée. Il possède alors une activité géométrique mais
pas d'activité décisionnelle détectée. Réciproquement, une faible
déformation peut traverser une séparatrice et devenir ponctuellement
décisionnelle.

La carte des propensions K/C doit donc être lue comme une propriété des
transitions observées sur un domaine, non comme une ontologie des
variables. Les taux empiriques mesurent la fréquence avec laquelle une
direction rencontre $K_\kappa$ au voisinage d'une bascule ; ils ne
constituent pas des probabilités universelles.

\section{Vérification dans le modèle canonique}
\label{annexe-J-verification}

Les observations numériques associées illustrent les quatre cas de la
proposition sans intervenir dans sa preuve.

\paragraph{Axe de pression $p$.}
Sur la plage déclarée, le compte et l'ordre des événements restent
constants alors que les dates dérivent. Les transitions d'issue
observées se produisent donc à l'intérieur de cellules hybrides : elles
sont de type C. Les micro-bandes d'appropriation correspondent à plusieurs
composantes du niveau de réponse dans une même géométrie cellulaire.

\paragraph{Durée $T_2$.}
À l'intérieur d'une marche de calendrier, modifier $T_2$ déplace l'état
au moment de la coupure sans recomposer la succession des événements.
Les transitions sont de type C. Aux bords des marches, le même axe peut
rencontrer $K_\kappa$.

\paragraph{Composition du préfixe.}
Les changements de composition étudiés antérieurement modifient la
succession des événements et produisent des transitions K sur les cas
mesurés. Il s'agit d'une propension empirique forte, non d'un théorème
sur tout paramètre de protocole.

\paragraph{État lent $(\theta,s)$.}
La majorité des transitions mesurées sont C, mais certaines recomposent
le calendrier et sont K. Ce contre-exemple suffit à invalider toute
dichotomie état/protocole.

\section{Ce que la proposition ne suppose pas}
\label{annexe-J-non-hypotheses}

La Proposition~\ref{prop:J-pli-couture} ne suppose pas :
\begin{itemize}
\item que $K_\kappa$ soit une variété ;
\item que $\Sigma$ soit globalement lisse ou connexe ;
\item que chaque recomposition change l'issue ;
\item que chaque bascule change le nombre d'événements ;
\item que les dates des événements soient constantes dans une cellule ;
\item que les catégories C et K soient exclusives aux points
singuliers ;
\item qu'un axe soit uniformément C ou K sur toute sa plage ;
\item qu'une classification obtenue à une résolution reste valide à une
autre résolution.
\end{itemize}

\section{Limites de portée}
\label{annexe-J-limites}

\paragraph{Signature insuffisante.}
Si $\kappa$ omet une donnée discrète qui modifie la factorisation de
l'exécution, une fausse cellule peut réunir deux mécanismes différents.
La signature doit être raffinée.

\paragraph{Signature excessivement fine.}
Si les dates exactes sont incluses dans $\kappa$, toute dérive continue
crée un changement de signature et détruit la distinction recherchée.
La signature doit rester combinatoire.

\paragraph{Non-unicité de l'exécution.}
En présence de solutions multiples ou de conflits événementiels non
résolus, $\kappa$ et $\widehat\Omega$ deviennent multivoques. La
proposition doit être remplacée par une version en correspondances.

\paragraph{Plateaux pré-binaires.}
Si $g_C=0$ sur un ouvert, le niveau zéro n'est pas une frontière de signe
ordinaire. Il faut définir la convention de lecture sur le plateau et
travailler avec la frontière topologique de l'ensemble d'issue.

\paragraph{Tangences.}
Un chemin peut toucher $\Sigma$ sans changer d'issue. L'appartenance à
$\gamma^{-1}(\Sigma)$ localise les candidats, non les bascules effectives.

\paragraph{Points mixtes.}
À l'intersection d'un pli cellulaire et de $K_\kappa$, attribuer un seul
mécanisme peut être arbitraire. La catégorie mixte est alors la seule
formulation stable.

\paragraph{Lectures historiques intégrales.}
Si la réponse lit l'histoire entière, la factorisation cellulaire doit
porter sur l'exécution segmentée complète. Le transport par préfixe
requiert la compatibilité au suffixe de l'Appendice~G.

\paragraph{Changement de loi.}
Une loi plastique peut modifier à la fois les cartes $E_C$, les niveaux
$g_C$ et le lieu $K_\kappa$. Comparer deux lois exige donc de distinguer
déplacement de frontière à signature fixe et recomposition de la
partition elle-même.

\section{Dépendances et ouverture vers la plasticité}
\label{annexe-J-dependances}

La chaîne logique est
\[
\begin{aligned}
&\text{Appendice G : segmentation et transport}\\
&\qquad\Longrightarrow \text{Appendices H--I : état suffisant}\\
&\qquad\Longrightarrow \text{partition par signatures}\\
&\qquad\Longrightarrow \text{Proposition~\ref{prop:J-pli-couture}}.
\end{aligned}
\]

L'annexe suivante ajoute une variable lente de loi $\gamma$. Deux
présentations pourront alors produire la même réponse en augmentant
l'état, tandis que deux lois différentes pourront déformer les plis
$\Sigma_C$ ou déplacer les coutures $K_\kappa$. La distinction établie
ici est donc nécessaire pour ne pas attribuer à la plasticité une
variation qui proviendrait seulement d'une recomposition du protocole.

\section{Table des notations}
\label{annexe-J-notations}

\begin{center}
\begin{tabular}{p{0.25\textwidth}p{0.62\textwidth}}
\toprule
Notation & Signification \\
\midrule
$D_{\mathrm c}$ & domaine causal des paramètres admissibles \\
$D_{\mathrm m}$ & domaine où le résumé déclaré est markovien \\
$D_{\mathrm r}$ & domaine réalisable par les histoires considérées \\
$\kappa$ & signature combinatoire de l'exécution \\
$C_k$ & intérieur relatif d'une classe de signature constante \\
$K_\kappa$ & lieu où la signature n'est pas localement constante \\
$E_C$ & carte d'état de coupure ou d'observation dans une cellule \\
$g_C$ & observable pré-binaire cellulaire \\
$\Sigma_C$ & partie de la frontière située dans une cellule \\
$\Sigma_K$ & partie de la frontière située sur le lieu de recomposition \\
C & transition à signature localement constante \\
K & transition rencontrant une recomposition de signature \\
\bottomrule
\end{tabular}
\end{center}

\section{Conclusion}
\label{annexe-J-conclusion}

La géométrie de la frontière ne se laisse pas réduire à l'opposition
entre paramètres d'état et paramètres de protocole. La distinction
stable porte sur l'exécution elle-même. Tant que sa signature
combinatoire demeure fixe, les paramètres déforment l'état parcouru et
plient la frontière à l'intérieur d'une cellule hybride. Lorsque la
signature change, la partition des exécutions se recompose et les
morceaux de frontière peuvent être recousus sur le lieu de changement.

Cette formulation est volontairement relative à une signature, une
résolution et un domaine réalisable. Elle conserve néanmoins un contenu
non tautologique : elle sépare les variations continues d'une exécution
de type fixe des changements de structure événementielle, et indique
exactement quelles hypothèses sont nécessaires pour passer d'une
observation numérique à une classification mathématique.
